\begin{table}[H]
\centering
\begin{tabularx}{\textwidth}{
  >{\centering\arraybackslash}X  % 符号列：占1/3宽度，居中
  >{\centering\arraybackslash}p{0.667\textwidth}  % 含义列：占2/3宽度，居中
  }
    \toprule[1.5pt]  % 顶部粗线
    \midrule
    符号 & 含义 \\
    \midrule
    $G$ & 剩余金额 \\

$N_t$ & 第$t$日所在区域编号 \\

$\Omega_s = \{N_s\}$ & 起点编号集合 \\

$\Omega_e = \{N_e\}$ & 终点编号集合 \\

$\Omega_v = \{N_{v_i} \mid i\in\mathbb{N}^*\}$ & 村庄编号集合 \\

$\Omega_h = \{N_{h_i} \mid i\in\mathbb{N}^*\}$ & 矿山编号集合 \\

$D$ & 天气（1-晴朗，2-高温，3-沙暴） \\

$X$ & 动向（1-移动，2-停留，3-挖矿） \\

$T = \max t$ & 到终点时的天数 \\

$n$ & 玩家总数 \\

$i\in\{1,2,\dots,n\}$ & 玩家编号 \\

$E$ & 地图邻接边集合 \\

$a_{uv}$ & 邻接矩阵元素，若区域 $u,v$ 相邻则 $a_{uv}=1$，否则 $a_{uv}=0$ \\

$Y_{i,t}^{uv}$ & 第 $i$ 名玩家在第 $t$ 日是否沿边 $(u,v)$ 移动的0-1变量 \\

$Z_{i,t}^{h}$ & 第 $i$ 名玩家在第 $t$ 日是否在矿山 $h\in\Omega_h$ 挖矿的0-1变量 \\

$Q_{i,t}^{v}$ & 第 $i$ 名玩家在第 $t$ 日是否在村庄 $v\in\Omega_v$ 购买资源的0-1变量 \\

$k_t^{uv}$ & 第 $t$ 日沿边 $(u,v)$ 同时移动的玩家数 \\

$k_t^{h}$ & 第 $t$ 日在矿山 $h$ 同时挖矿的玩家数 \\

$k_t^{v}$ & 第 $t$ 日在村庄 $v$ 同时购买资源的玩家数 \\

$\delta$ & 联盟稳定性指标 \\

$\mathcal P_i$ & 由前两问路径模型生成的第 $i$ 名玩家候选可行路径集 \\

$\pi_t(d'\mid d)$ & 天气状态从 $d$ 转移到 $d'$ 的条件概率 \\

$V_t(\cdot)$ & 第 $t$ 日滚动决策价值函数 \\
    \midrule

    \bottomrule[1.5pt]  % 底部粗
\end{tabularx}
\end{table}